\documentclass[11pt]{article}
\usepackage[utf8]{inputenc}
\usepackage[T1]{fontenc}
\usepackage{geometry}
\geometry{a4paper, margin=1in}
\usepackage{enumitem}
\usepackage{titlesec}
\usepackage{fancyhdr}
\usepackage{lastpage}

% Header/Footer
\pagestyle{fancy}
\fancyhf{}
\rhead{Prompt Engineering - Final Exam}
\lhead{Ankara Medipol University}
\rfoot{Page \thepage\ of \pageref{LastPage}}

\title{\textbf{Prompt Engineering}\\\large Final Examination\\Spring 2025--2026}
\author{Instructor: Mehmet Ali Akyol, PhD\\Ankara Medipol University}
\date{}

\begin{document}

\maketitle

\noindent\textbf{Student Name:} \underline{\hspace{6cm}} \hfill \textbf{Student ID:} \underline{\hspace{4cm}}

\vspace{0.5cm}
\noindent\textbf{Instructions:}
\begin{itemize}[noitemsep]
    \item This exam consists of 40 multiple-choice questions.
    \item Each question is worth 2.5 points for a total of 100 points.
    \item Select the \textbf{best} answer for each question.
    \item You have 90 minutes to complete this exam.
    \item No electronic devices or notes are permitted.
\end{itemize}

\vspace{0.5cm}
\hrule
\vspace{0.5cm}

\begin{enumerate}

\item A \textbf{prompt} in the context of AI is best defined as:
\begin{enumerate}[label=\Alph*)]
    \item The internal memory of an LLM
    \item The training data used to build the model
    \item The hardware that runs the AI system
    \item The input message you give to an AI system to elicit a response
\end{enumerate}

\item Which of the following is NOT one of the core elements that good prompts provide?
\begin{enumerate}[label=\Alph*)]
    \item Purpose
    \item Context
    \item Constraints
    \item Source code
\end{enumerate}

\item What is \textbf{tokenization} in the context of LLMs?
\begin{enumerate}[label=\Alph*)]
    \item The process of encrypting user data
    \item Generating authentication tokens for API access
    \item Converting images into text descriptions
    \item Breaking text into smaller chunks (tokens) that the model can process
\end{enumerate}

\item Modern LLMs like GPT-4 use subword tokenization primarily because:
\begin{enumerate}[label=\Alph*)]
    \item It makes the model run faster
    \item It reduces the cost of training
    \item It improves image recognition
    \item It can handle new or rare words by combining familiar pieces
\end{enumerate}

\item What are \textbf{embeddings} in LLM architecture?
\begin{enumerate}[label=\Alph*)]
    \item Physical components of the computer running the model
    \item The user interface elements of AI applications
    \item Error logs generated during model training
    \item Numerical representations of tokens that capture semantic meaning
\end{enumerate}

\item The 2017 paper ``Attention Is All You Need'' introduced which key innovation?
\begin{enumerate}[label=\Alph*)]
    \item Self-attention mechanism (Transformer architecture)
    \item Tokenization
    \item Recurrent neural networks
    \item Convolutional neural networks
\end{enumerate}

\item \textbf{Hallucination} in LLMs refers to:
\begin{enumerate}[label=\Alph*)]
    \item The model refusing to answer questions
    \item The model asking clarifying questions
    \item The model generating plausible but incorrect or fabricated information
    \item The model running out of memory
\end{enumerate}

\item Which statement about LLM capabilities is TRUE?
\begin{enumerate}[label=\Alph*)]
    \item LLMs have true understanding and consciousness
    \item LLMs have unlimited context windows
    \item LLMs are statistical pattern learners, not conscious reasoners
    \item LLMs never produce biased outputs
\end{enumerate}

\item In the anatomy of a good prompt, which element specifies headings, JSON schema, or table columns?
\begin{enumerate}[label=\Alph*)]
    \item Task
    \item Context
    \item Constraints
    \item Output Specification
\end{enumerate}

\item \textbf{Zero-shot prompting} is best described as:
\begin{enumerate}[label=\Alph*)]
    \item Giving just the task without any examples
    \item Providing multiple examples before the task
    \item Using images along with text
    \item Chaining multiple prompts together
\end{enumerate}

\item \textbf{Few-shot prompting} is most useful when:
\begin{enumerate}[label=\Alph*)]
    \item You want to save tokens
    \item The task is extremely simple
    \item You need to steer style or format with specific examples
    \item You want to reduce response time
\end{enumerate}

\item What is \textbf{role prompting}?
\begin{enumerate}[label=\Alph*)]
    \item Asking the model to play a video game
    \item Defining database roles
    \item Setting a persona, audience, and tone for the model's response
    \item Assigning user permissions to the model
\end{enumerate}

\item Which is an example of a ``prompt anti-pattern''?
\begin{enumerate}[label=\Alph*)]
    \item Specifying output format clearly
    \item Providing relevant context
    \item Setting clear constraints
    \item Including conflicting directives like ``be brief'' and ``include all details''
\end{enumerate}

\item In Retrieval-Augmented Generation (RAG), what is the primary purpose?
\begin{enumerate}[label=\Alph*)]
    \item To speed up model inference
    \item To reduce the model size
    \item To ground answers with provided excerpts or references
    \item To train the model on new data
\end{enumerate}

\item The \textbf{CO-STAR} framework stands for:
\begin{enumerate}[label=\Alph*)]
    \item Code, Output, Style, Testing, Analysis, Review
    \item Context, Objective, Style, Tone, Audience, Response
    \item Constraint, Optimization, Structure, Target, Action, Result
    \item Chain, Order, Sequence, Task, Answer, Repeat
\end{enumerate}

\item The \textbf{IDEA Loop} framework for prompt iteration consists of:
\begin{enumerate}[label=\Alph*)]
    \item Implement, Debug, Edit, Apply
    \item Input, Design, Execute, Analyze
    \item Initiate, Develop, Enhance, Archive
    \item Identify, Draft, Evaluate, Adjust
\end{enumerate}

\item \textbf{Meta-prompting} refers to:
\begin{enumerate}[label=\Alph*)]
    \item Using prompts to have an LLM generate or refine other prompts
    \item Prompting about metadata
    \item Using multiple models simultaneously
    \item Prompting in multiple languages
\end{enumerate}

\item In the context of tool use and function calling, the LLM typically:
\begin{enumerate}[label=\Alph*)]
    \item Directly executes code on the user's computer
    \item Sends requests to arbitrary websites
    \item Generates a structured object (e.g., JSON) specifying the tool and arguments
    \item Modifies its own training data
\end{enumerate}

\item Which dimension of linguistic register controls terminology density?
\begin{enumerate}[label=\Alph*)]
    \item Tenor (Relationship)
    \item Mode (Channel)
    \item Field (Topic Domain)
    \item Rhythm
\end{enumerate}

\item ``CEFR B1'' in prompt specifications refers to:
\begin{enumerate}[label=\Alph*)]
    \item A programming language version
    \item A model architecture type
    \item A security classification
    \item A reading/language proficiency level
\end{enumerate}

\item \textbf{Brand Voice Fingerprinting} involves:
\begin{enumerate}[label=\Alph*)]
    \item Creating biometric authentication
    \item Recording audio samples of brand representatives
    \item Extracting measurable voice attributes from existing content to ensure AI outputs match brand style
    \item Generating unique identifiers for users
\end{enumerate}

\item The \textbf{Critique Loop} in prompt refinement typically follows which sequence?
\begin{enumerate}[label=\Alph*)]
    \item Generate $\to$ Deploy $\to$ Monitor
    \item Design $\to$ Test $\to$ Archive
    \item Plan $\to$ Execute $\to$ Report
    \item Generate $\to$ Critique $\to$ Revise
\end{enumerate}

\item ``Vibe Coding'' refers to:
\begin{enumerate}[label=\Alph*)]
    \item Writing code by describing desired behavior to an LLM and iterating based on output's look and feel
    \item Writing code based on musical rhythms
    \item A formal software development methodology
    \item Debugging code using vibration sensors
\end{enumerate}

\item What is the main advantage of using ``LLM-as-a-Judge'' for evaluation?
\begin{enumerate}[label=\Alph*)]
    \item It is always 100\% accurate
    \item It eliminates the need for any human review
    \item It enables scalable evaluation of thousands of outputs quickly
    \item It reduces model training costs
\end{enumerate}

\item The ``3 H's'' framework for AI output evaluation includes Helpful, Honest, and:
\begin{enumerate}[label=\Alph*)]
    \item Historical
    \item Hypothetical
    \item Harmless
    \item Hierarchical
\end{enumerate}

\item A \textbf{Golden Dataset} in AI evaluation is:
\begin{enumerate}[label=\Alph*)]
    \item A curated set of inputs with ``perfect'' outputs used as benchmarks
    \item The largest available training dataset
    \item Data stored in a premium database
    \item User feedback collected over time
\end{enumerate}

\item \textbf{Verbosity Bias} in LLM evaluation refers to:
\begin{enumerate}[label=\Alph*)]
    \item Models producing too few tokens
    \item Preference for technical jargon
    \item The tendency to rate longer answers as better, even if they contain fluff
    \item Bias toward English-language responses
\end{enumerate}

\item In Vision-Language Models (VLMs), images are converted into:
\begin{enumerate}[label=\Alph*)]
    \item Audio signals
    \item Database entries
    \item Encrypted files
    \item Patches that become ``visual tokens''
\end{enumerate}

\item \textbf{Set-of-Mark (SoM) Prompting} involves:
\begin{enumerate}[label=\Alph*)]
    \item Overlaying numbered tags or bounding boxes on images to improve grounding
    \item Grading student assignments
    \item Creating watermarks on AI-generated images
    \item Setting bookmarks in long documents
\end{enumerate}

\item What is \textbf{Visual Hallucination}?
\begin{enumerate}[label=\Alph*)]
    \item The model refusing to process images
    \item Image compression artifacts
    \item The model describing objects or details that aren't present in the image
    \item Creating optical illusions
\end{enumerate}

\item The \textbf{MECE Principle} stands for:
\begin{enumerate}[label=\Alph*)]
    \item Mutually Exclusive, Collectively Exhaustive
    \item Maximum Efficiency, Controlled Execution
    \item Model Evaluation, Context Enhancement
    \item Multiple Examples, Chain Execution
\end{enumerate}

\item \textbf{Tree of Thoughts (ToT)} differs from Chain-of-Thought by:
\begin{enumerate}[label=\Alph*)]
    \item Using fewer tokens
    \item Being faster to execute
    \item Requiring no examples
    \item Exploring multiple reasoning branches and allowing backtracking
\end{enumerate}

\item The \textbf{ReAct} framework combines:
\begin{enumerate}[label=\Alph*)]
    \item Recalling and Testing
    \item Reading and Acting
    \item Reasoning and Acting
    \item Responding and Chatting
\end{enumerate}

\item In \textbf{Self-Consistency} prompting, you:
\begin{enumerate}[label=\Alph*)]
    \item Ask the model to repeat the same answer multiple times
    \item Force the model to be consistent with previous conversations
    \item Use the same prompt for different models
    \item Generate multiple independent answers using different reasoning paths and pick the most consistent
\end{enumerate}

\item The \textbf{SOAP Note} format (Subjective, Objective, Assessment, Plan) is commonly used in which domain?
\begin{enumerate}[label=\Alph*)]
    \item Legal
    \item Finance
    \item Healthcare
    \item Engineering
\end{enumerate}

\item The \textbf{IRAC} framework (Issue, Rule, Analysis, Conclusion) is specifically designed for:
\begin{enumerate}[label=\Alph*)]
    \item Medical diagnosis
    \item Financial reporting
    \item Software development
    \item Legal reasoning and analysis
\end{enumerate}

\item A \textbf{Knowledge Pack} in domain-specific prompting typically includes:
\begin{enumerate}[label=\Alph*)]
    \item Definitions, acronyms, and style guides specific to the domain
    \item Only training data
    \item Password credentials
    \item Model weights
\end{enumerate}

\item The term \textbf{Automation Bias} refers to:
\begin{enumerate}[label=\Alph*)]
    \item Bias in robotic systems
    \item Preference for automated testing
    \item The tendency for humans to over-rely on AI outputs without critical verification
    \item AI systems preferring automated tasks
\end{enumerate}

\item What is \textbf{Prompt Injection}?
\begin{enumerate}[label=\Alph*)]
    \item Adding more tokens to a prompt
    \item Injecting code into the model's training
    \item A medical procedure
    \item Crafting inputs to deliberately bypass AI safety protocols
\end{enumerate}

\item \textbf{Agentic AI} differs from traditional chatbots by being:
\begin{enumerate}[label=\Alph*)]
    \item Unable to use external tools
    \item Text-only without multimodal capabilities
    \item Proactive, multi-step, stateful, and capable of autonomous tool usage
    \item Slower and less accurate
\end{enumerate}

\end{enumerate}

\vspace{1cm}
\hrule
\vspace{0.5cm}

\begin{center}
\textbf{--- END OF EXAM ---}
\end{center}

\newpage
%=============================================================================
\section*{Answer Key (For Instructor Use Only)}
%=============================================================================

\begin{center}
\begin{tabular}{|c|c||c|c||c|c||c|c|}
\hline
\textbf{Q} & \textbf{Ans} & \textbf{Q} & \textbf{Ans} & \textbf{Q} & \textbf{Ans} & \textbf{Q} & \textbf{Ans} \\
\hline
1 & D & 11 & C & 21 & C & 31 & A \\
2 & D & 12 & C & 22 & D & 32 & D \\
3 & D & 13 & D & 23 & A & 33 & C \\
4 & D & 14 & C & 24 & C & 34 & D \\
5 & D & 15 & B & 25 & C & 35 & C \\
6 & A & 16 & D & 26 & A & 36 & D \\
7 & C & 17 & A & 27 & C & 37 & A \\
8 & C & 18 & C & 28 & D & 38 & C \\
9 & D & 19 & C & 29 & A & 39 & D \\
10 & A & 20 & D & 30 & C & 40 & C \\
\hline
\end{tabular}
\end{center}

\vspace{1cm}

\section*{Detailed Answer Key with Explanations}

\begin{enumerate}
    \item \textbf{D} - A prompt is the input message given to an AI system.
    \item \textbf{D} - Source code is not a core element; purpose, context, and constraints are.
    \item \textbf{D} - Tokenization breaks text into processable chunks.
    \item \textbf{D} - Subword tokenization handles new/rare words flexibly.
    \item \textbf{D} - Embeddings are numerical vectors capturing semantic meaning.
    \item \textbf{A} - The Transformer architecture introduced self-attention.
    \item \textbf{C} - Hallucination is generating plausible but incorrect information.
    \item \textbf{C} - LLMs are pattern learners, not conscious reasoners.
    \item \textbf{D} - Output Specification defines format like JSON, tables.
    \item \textbf{A} - Zero-shot means no examples, just the task.
    \item \textbf{C} - Few-shot steers style/format with examples.
    \item \textbf{C} - Role prompting sets persona, audience, tone.
    \item \textbf{D} - Conflicting directives are a common anti-pattern.
    \item \textbf{C} - RAG grounds answers with provided context.
    \item \textbf{B} - CO-STAR: Context, Objective, Style, Tone, Audience, Response.
    \item \textbf{D} - IDEA Loop: Identify, Draft, Evaluate, Adjust.
    \item \textbf{A} - Meta-prompting uses prompts to generate/refine prompts.
    \item \textbf{C} - Tool use generates structured JSON for function calls.
    \item \textbf{C} - Field (Topic Domain) controls terminology density.
    \item \textbf{D} - CEFR B1 is a language proficiency level.
    \item \textbf{C} - Brand Voice Fingerprinting extracts measurable attributes.
    \item \textbf{D} - Critique Loop: Generate $\to$ Critique $\to$ Revise.
    \item \textbf{A} - Vibe Coding is iterating on LLM output based on feel.
    \item \textbf{C} - LLM-as-Judge enables scalable evaluation.
    \item \textbf{C} - The 3 H's: Helpful, Honest, Harmless.
    \item \textbf{A} - Golden Dataset is a curated benchmark set.
    \item \textbf{C} - Verbosity Bias favors longer answers.
    \item \textbf{D} - VLMs convert images into visual tokens.
    \item \textbf{A} - Set-of-Mark overlays numbered markers on images.
    \item \textbf{C} - Visual Hallucination describes non-existent details.
    \item \textbf{A} - MECE: Mutually Exclusive, Collectively Exhaustive.
    \item \textbf{D} - ToT explores multiple branches with backtracking.
    \item \textbf{C} - ReAct combines Reasoning and Acting.
    \item \textbf{D} - Self-Consistency generates multiple answers to find consensus.
    \item \textbf{C} - SOAP Notes are used in Healthcare.
    \item \textbf{D} - IRAC is for Legal reasoning.
    \item \textbf{A} - Knowledge Packs include definitions, acronyms, style guides.
    \item \textbf{C} - Automation Bias is over-reliance on AI outputs.
    \item \textbf{D} - Prompt Injection bypasses AI safety protocols.
    \item \textbf{C} - Agentic AI is proactive, multi-step, stateful.
\end{enumerate}

\vspace{1cm}

\section*{Topic Coverage by Question}

\begin{itemize}
    \item Foundations: Q1--Q8
    \item Prompt Design Principles: Q9--Q14
    \item Patterns and Frameworks: Q15--Q18
    \item Style, Tone, and Format Control: Q19--Q21
    \item Iterative Refinement: Q22--Q24
    \item Evaluation Methods: Q25--Q27
    \item Multimodal Prompting: Q28--Q30
    \item Complex Problem Solving: Q31--Q34
    \item Domain-Specific Applications: Q35--Q37
    \item Ethics, Responsible AI \& Emerging Trends: Q38--Q40
\end{itemize}

\end{document}
